% Generated draft: requires mathematical and source review.
\documentclass[11pt]{article}
\usepackage[margin=25mm]{geometry}
\usepackage{fontspec}
\setmainfont{Latin Modern Roman}
\usepackage{amsmath,amssymb,mathtools,unicode-math}
\setmathfont{Latin Modern Math}
\usepackage{enumitem,needspace,xurl,hyperref}
\setlength{\parindent}{0pt}
\setlength{\parskip}{5pt}
\setlength{\emergencystretch}{3em}
\newcommand{\N}{\mathbb{N}}
\newcommand{\Z}{\mathbb{Z}}
\newcommand{\Q}{\mathbb{Q}}
\newcommand{\R}{\mathbb{R}}
\newcommand{\C}{\mathbb{C}}
\newcommand{\F}{\mathbb{F}}
\newcommand{\A}{\mathbb{A}}
\newcommand{\PP}{\mathbb{P}}
\newcommand{\E}{\mathbb{E}}
\newcommand{\eps}{\varepsilon}
\newcommand{\dd}{\,\mathrm{d}}
\begin{document}
\section*{UnsolvedMath ID: 10400055}
\textbf{Source problem code:} catalogue entry
\subsection{Definitions and mathematical statement}
Let \(Z(K)\) denote the Kontsevich integral of a knot \(K\), with the normalization, framing, coefficient field and associator conventions of the source still to be specified. The notation \(\widehat{\mathcal A}(S^1)\) ordinarily denotes a degree-completed algebra of Jacobi diagrams on an oriented circle skeleton, modulo the relevant diagram relations. The suffix \({\rm conn}\) denotes its connected/primitive part in the source's convention. The notation \(\mathcal B_{\rm conn}\) refers to the corresponding connected diagram space without the circle skeleton, but the record does not state its completion or the map used to transport \(Z\) into it.

For a complete filtered algebra with product \(*\), unit \(1\), and \(x-1\) of positive filtration, the formal logarithm is
\[
 \log_*(x)=\sum_{m\ge1}\frac{(-1)^{m+1}}{m}(x-1)^{*m},
\]
with convergence in the filtration; the coefficients require a suitable characteristic-zero coefficient system. In the circle-diagram algebra this gives the source's \(\log Z(K)\); in the skeleton-free algebra the product indicated by \(\log_{\sqcup}\) is disjoint union. This explains the formal operation without identifying two potentially different algebra structures.

Problem 3.5 asks for a characterization of the realizable subsets
\[
 \{\log Z(K):K\text{ is a knot in the source's class}\}
       \subset\widehat{\mathcal A}(S^1)_{\rm conn}
\]
and, in the alternative formulation,
\[
 \{\log_{\sqcup} Z(K):K\text{ is a knot in the source's class}\}
       \subset\mathcal B_{\rm conn}.
\]
The desired characterization is of which diagram series actually arise from individual knots, not merely their linear span. The source supplies neither a proposed criterion nor enough diagram/normalization conventions to declare these two displayed formulations equivalent. Relations, framing quotients, completion on the \(\mathcal B\) side, and the precise skeleton-removal map must be checked before acceptance.
\subsection{Short English statement}
Characterize the connected diagram series that occur as logarithms of Kontsevich integrals of actual knots, either in the circle-skeleton algebra or using the disjoint-union logarithm in the skeleton-free space. The source record leaves the normalization, completed spaces, and relation between the two formulations implicit.
\subsection{Sources}
\begin{itemize}
\item Local UnsolvedMath registry, permanent ID \(10400055\), exact source code \texttt{catalogue entry}, Problem 3.5. No individual proposer or proposed year is supplied.
\item Ohtsuki, \emph{Problems on invariants of knots and 3-manifolds} (2002), Problem 3.5, PDF page 53: \url{https://msp.org/gtm/2002/04/gtm-2002-04-024s.pdf}. Printed bold headings distinguish 231 numbered items from surrounding references.
\item Extraction pdf-font-structure, accessed 2026-07-29. Source evidence NEEDS\_REVIEW presents an open problem and requires a current primary-literature status audit. Supplied research classification OPEN-TRIAGE identifies the question with the image of the Kontsevich integral; no further reference or resolution is supplied.
\item Dataset attribution: Ulam AI / UnsolvedMath Contributors, CC BY 4.0. Source rights NEEDS\_REVIEW: publicly accessible author/publisher copy, redistribution terms requiring review.
\end{itemize}
\end{document}
